\section{Implementation Details}

\subsection{Software Stack}
The implementation combines several technology choices selected for portability, safety, and compatibility with the target operational model.
\begin{itemize}
    \item \textbf{Rust services.} The API server, gateway, and controllers are implemented in Rust, leveraging the language's memory-safety guarantees and ecosystem support for asynchronous networked services~\cite{rust}. This choice is particularly appropriate for long-running control-plane components where concurrency and resilience matter.
    \item \textbf{Kubernetes-native control logic.} Resource interaction follows controller and status-patching patterns consistent with Kubernetes operational practice~\cite{kubernetes,kubernetes_operators}. This allows RETROSPECT to treat embedded endpoints as first-class managed resources without pretending that they are standard cluster nodes.
    \item \textbf{Zephyr firmware.} Device-side communication and execution support are grounded in Zephyr, which provides a mature embedded software substrate with networking and TLS capabilities~\cite{zephyr}.
    \item \textbf{WAMR runtime.} Edge-side application execution uses WAMR, enabling lightweight WASM deployment semantics on constrained targets~\cite{wamr}.
\end{itemize}

The choice of Rust on the control-plane side and Zephyr/WAMR on the device side is not accidental. Together they support a layered trust argument: memory-safe control services reduce infrastructure fragility, while a lightweight RTOS and runtime pair keep the device footprint consistent with embedded requirements.

\subsection{Deployment Topology}
The reference deployment runs on K3s~\cite{k3s} with namespace-scoped resources and locally built images. This topology is deliberately modest: it preserves real Kubernetes semantics while remaining accessible on a single-node Linux environment suitable for development, experimentation, and artifact evaluation. The deployment includes API services, dashboard, gateway/controller components, CRDs, and the supporting services required to exercise enrollment and status propagation.

From a reproducibility standpoint, the use of K3s is important because it lowers the infrastructure barrier while retaining the reconciliation behavior and API machinery needed to make the results representative of cluster-native operation. The system can therefore be evaluated without a large cluster but without reducing the experiment to a mock or non-Kubernetes approximation.

\subsection{Controller Responsibilities}
The implementation distributes responsibilities across specialized controllers rather than concentrating all logic into a single service. Device-oriented reconciliation manages attachment state and gateway resolution. Application-oriented reconciliation interprets desired deployment intent and maps it to gateway actions. Gateway-oriented reconciliation maintains metadata needed for path resolution and status interpretation. This decomposition aligns with the CRD model and reduces the risk that unrelated concerns become coupled inside one control loop.

\subsection{Reconciliation Hardening}
During validation, a critical defect emerged: a connected device could be incorrectly forced into a \texttt{Disconnected} phase due to transient checks and empty gateway endpoint propagation. The defect was not in the enrollment protocol itself, which completed correctly, but in the policy used to interpret the resulting evidence. The corrective action involved three changes:
\begin{itemize}
    \item persisting a valid default gateway endpoint when the status path temporarily lacked one;
    \item avoiding aggressive forced disconnection when transport visibility briefly diverged from the latest known heartbeat evidence;
    \item preserving state continuity until the next reconciliation cycle could confirm whether the condition was real or transient.
\end{itemize}

This hardening is technically significant because it highlights a common risk in cyber-physical orchestration systems: transport-level truth and control-plane truth can diverge temporarily, and naive reconciliation can transform a transient mismatch into a persistent semantic error. RETROSPECT therefore treats reconciliation policy as a first-class design element rather than an implementation detail.

\subsection{Operational Lessons from the Implementation}
A practical outcome of the implementation effort is that end-to-end correctness depends on the integration details surrounding the core logic. Image availability, cluster runtime semantics, protocol framing, controller policies, and status-patching behavior all influence whether the platform works as intended. For this reason, the paper documents not only the nominal architecture but also the operational decisions required to produce a working deployment. This level of detail is necessary for a submission that aims to be technically credible and reproducible rather than merely conceptual.
